
\textbf{Актуальность} связана с необходимостью выравнивания LLM в более сложных стратегических и этических средах, включая мульти-агентные системы.

% {\progress}
% Этот раздел должен быть отдельным структурным элементом по
% ГОСТ, но он, как правило, включается в описание актуальности
% темы. Нужен он отдельным структурынм элемементом или нет ---
% смотрите другие диссертации вашего совета, скорее всего не нужен.

{\aim} данной работы является исследование влияния эмпатии в больших
языковых моделях при их использовании в процессе принятия решений в
теории игр, разработка экспериментальной методологии для анализа влияния
эмоций в играх с торгом, подбор проприетарных и свободных моделей
искусственного интеллекта, проведение экспериментов в многоязычных
средах с разными эмоциональными состояниями, анализ результатов.

Для~достижения поставленной цели необходимо было решить следующие {\tasks}:
\begin{enumerate}[beginpenalty=10000] % https://tex.stackexchange.com/a/476052/104425
  \item Рассмотреть теоретические аспекты влияния эмпатии в контексте теории игр.
  \item Разработать фреймворк для проведения экспериментов для оценки влияния эмоций в больших языковых моделях на примере игр <<Диктатор>>, <<Дилемма узника>> и <<Общественное благо>>.
  \item Провести эксперименты, связанные с эмоциональным воздействием, в исследуемых играх.
  \item Выявить модели, наиболее соответствующие выравниванию с эмоциональным поведением человека.
\end{enumerate}


{\novelty}
%\begin{enumerate}[beginpenalty=10000] % https://tex.stackexchange.com/a/476052/104425
Была проведена оценка влияний эмпатии на этическое и эмоциональное поведение LLM в контексте теории игр.
%\end{enumerate}

{\influence} Практическая значимость эмпатии для выравнивания в крупных языковых моделях заключается в её возможности увеличить эффективность в реальных жизненных ситуациях, не имеющих одногозначного решения. 

{\methods} В работе применялись методы теории игр, статистический анализ данных, методы управления большими языковыми моделями, общенаучные методы анализ и синтез.

%{\defpositions}
%\begin{enumerate}[beginpenalty=10000] % https://tex.stackexchange.com/a/476052/104425
%  \item Первое положение
%  \item Второе положение
%  \item Третье положение
%  \item Четвертое положение
%\end{enumerate}
%В папке Documents можно ознакомиться с решением совета из Томского~ГУ
%(в~файле \verb+Def_positions.pdf+), где обоснованно даются рекомендации
%по~формулировкам защищаемых положений.

%{\reliability} полученных результатов обеспечивается \ldots \ Результаты находятся в соответствии с результатами, полученными другими авторами.
